\documentclass[10pt,landscape]{article}

\usepackage[margin=0.32in]{geometry}
\usepackage[T1]{fontenc}
\usepackage[utf8]{inputenc}
\usepackage{lmodern}
\usepackage{textcomp}
\usepackage{microtype}
\usepackage[table,dvipsnames]{xcolor}
\usepackage{array}
\usepackage{booktabs}
\usepackage{enumitem}
\usepackage{hyperref}
\usepackage{listings}
\usepackage{multicol}
\usepackage{tabularx}
\usepackage{titlesec}

\pagestyle{empty}
\setlength{\parindent}{0pt}
\setlength{\parskip}{1.5pt}
\setlength{\columnsep}{13pt}
\setlength{\columnseprule}{0.2pt}
\setcounter{secnumdepth}{0}

\definecolor{ToshBlue}{HTML}{0F4C81}
\definecolor{ToshInk}{HTML}{1F2937}
\definecolor{ToshMuted}{HTML}{5B6470}
\definecolor{ToshPanel}{HTML}{F4F7FB}
\definecolor{ToshRule}{HTML}{CBD5E1}
\definecolor{ToshCodeBg}{HTML}{F8FAFC}

\color{ToshInk}

\titleformat{\section}
  {\large\bfseries\color{ToshBlue}}
  {}
  {0pt}
  {}
\titlespacing*{\section}{0pt}{0.4em}{0.12em}

\setlist[itemize]{leftmargin=1.05em,itemsep=0.5pt,topsep=1.5pt,parsep=0pt,partopsep=0pt}

\lstdefinestyle{tosh}{
  basicstyle=\ttfamily\scriptsize,
  backgroundcolor=\color{ToshCodeBg},
  frame=single,
  rulecolor=\color{ToshRule},
  framerule=0.4pt,
  xleftmargin=0pt,
  xrightmargin=0pt,
  aboveskip=0.18em,
  belowskip=0.24em,
  columns=fullflexible,
  keepspaces=true,
  showstringspaces=false,
  breaklines=true,
  literate={°}{{\textdegree}}1
           {μ}{{\ensuremath{\mu}}}1
           {Ω}{{\ensuremath{\Omega}}}1
}
\lstset{style=tosh}

\newcommand{\sheettitle}[3]{
  {\LARGE\bfseries\color{ToshBlue} ToSh Cheat Sheet: #1\par}
  \vspace{0.08em}
  {\small\color{ToshMuted} #2 \hfill #3\par}
  \vspace{0.18em}
  {\color{ToshRule}\rule{\linewidth}{0.8pt}}
  \vspace{0.12em}
}

\newcommand{\note}[1]{
  \vspace{0.12em}
  \fcolorbox{ToshRule}{ToshPanel}{%
    \parbox{\dimexpr\linewidth-2\fboxsep-2\fboxrule\relax}{\footnotesize #1}%
  }
  \vspace{0.12em}
}

\newcommand{\tighttable}[2]{
  {\footnotesize
  \renewcommand{\arraystretch}{1.08}
  \begin{tabularx}{\linewidth}{@{}#1@{}}
  #2
  \end{tabularx}}
}


\begin{document}
\footnotesize
\sheettitle{Collections And Pipelines}{Literals, typed containers, projection, grouping, and summaries}{spec + examples + tests}

\begin{multicols*}{3}

\section{Collection Literals}
\begin{lstlisting}
# arrays
[1, 2, 3]
["alpha", "beta"]

# record-like dynamic objects
{| Name = "Komrad", Uid = 1000 |}

# dictionaries with arbitrary keys
{% "name" => "Komrad", "uid" => 1000 %}

# sets
{: "alpha", "beta", "alpha" :}

# tuples
(42, "hello")
\end{lstlisting}

\tighttable{>{\ttfamily\raggedright\arraybackslash}p{0.26\linewidth}X}{
\toprule
Literal & Runtime shape \\
\midrule
\texttt{[...]} & array \\
\texttt{\{| Name = Value |\}} & dynamic record \\
\texttt{\{\% Key => Value \%\}} & dictionary \\
\texttt{\{: ... :\}} & set with de-duplication \\
\texttt{(..., ...)} & tuple with \texttt{Item1}, \texttt{Item2}, ... \\
\bottomrule
}

\section{Constructed Collections}
\begin{lstlisting}
new list<int>(10, 20, 30)
new array<int>(1, 2, 3)
new dict<string, int>(Alice, 95, Bob, 87)
new map({| Theme = "dark", TabSize = 4 |})
new hashtable(HOME, "/home/me", SHELL, "/usr/bin/tosh")
new record(Name, "Toast", Uid, 1000)
new tuple(alpha, 42)
\end{lstlisting}

\begin{itemize}
\item \texttt{list<T>} is growable and works well for accumulation.
\item \texttt{dict<K, V>} gives a strongly typed dictionary.
\item \texttt{map(...)} and record literals are handy for dynamic data.
\item \texttt{record} builds an anonymous dynamic record; \texttt{table} and \texttt{dynamicrecord} are aliases.
\end{itemize}

\section{Basic Access}
\begin{lstlisting}
var words = ["hello", "world"]
echo $words[0]
echo $words.Length

var pair = (42, "hello")
echo $pair.Item2
\end{lstlisting}

\begin{lstlisting}
var person = {| Name = "Komrad", Uid = 1000 |}
echo $person.Name
get-props $person
has-prop $person Name
\end{lstlisting}

\columnbreak

\section{Pipeline Core}
\tighttable{>{\ttfamily\raggedright\arraybackslash}p{0.27\linewidth}X}{
\toprule
Command & Use \\
\midrule
\texttt{where} & filter rows with a predicate \\
\texttt{get} / \texttt{pick} / \texttt{select} & project members or indices \\
\texttt{each} & run a block per item \\
\texttt{map} & transform each item \\
\texttt{flat-map} & transform and flatten \\
\texttt{flatten} & expand one nesting level \\
\texttt{sort}, \texttt{reverse} & reorder results \\
\texttt{first}, \texttt{last}, \texttt{skip} & slice a stream \\
\bottomrule
}

\begin{lstlisting}
ls -la | where _.Type == file | get { Name, Size }
echo "Hello" "World" | each { _.ToLower() }
echo [1, 2, 3] | flatten | map { _ * 2 }
\end{lstlisting}

\note{Inside pipeline predicates, \texttt{\_} is the current item. In many places a block like \texttt{\{ \_.Size > 1kb \}} is easier to read than a named lambda.}

\section{Grouping And Stats}
\begin{lstlisting}
ls -la | group-by Type
ls -la | frequencies
ls -la | distinct
\end{lstlisting}

\begin{lstlisting}
seq 5 | summarize --sum --avg --min --max --count
df | summarize _.Used
ps | summarize --avg WorkingSet --max WorkingSet
\end{lstlisting}

\tighttable{>{\ttfamily\raggedright\arraybackslash}p{0.28\linewidth}X}{
\toprule
Command & Best for \\
\midrule
\texttt{group-by} & collect rows by key \\
\texttt{group-while} & consecutive runs \\
\texttt{frequencies} & histogram-like counts \\
\texttt{summary} / \texttt{summarize} & aggregates over scalars or columns \\
\bottomrule
}

\section{Structured Text}
\begin{lstlisting}
echo "[{\"name\":\"alpha\",\"score\":95}]" | from json | flatten
echo "a,b\n1,2" | from csv
echo "toast" | hash sha256
echo "the quick brown fox" | grep "quick"
\end{lstlisting}

\begin{lstlisting}
$items | to json
$rows | to csv
$rows | template "{{ Name }} => {{ Quantity }}"
\end{lstlisting}

\columnbreak

\section{Collection Patterns}
\begin{lstlisting}
var items = new list<string>("bread", "coffee", "soap")
$items | where _.Length > 4 | sort | each { echo _ }
\end{lstlisting}

\begin{lstlisting}
var report = new list()
ls -la | where _.Type == file | each {
    $report.Add({| Name = _.Name, Size = _.Size |}) | ignore
}
$report | sort Size | reverse
\end{lstlisting}

\section{Transform Recipes}
\begin{lstlisting}
# top 5 files by size
ls -la | where _.Type == file | sort Size | reverse | first 5

# names only
ls -la | get Name

# files grouped by extension
ls -la | where _.Type == file | group-by Extension
\end{lstlisting}

\begin{lstlisting}
# flatten JSON rows then rank them
echo "[{\"name\":\"a\",\"score\":1},{\"name\":\"b\",\"score\":2}]" |
from json | flatten | sort score | reverse | get { name, score }
\end{lstlisting}

\section{Small But Handy}
\tighttable{>{\ttfamily\raggedright\arraybackslash}p{0.29\linewidth}X}{
\toprule
Tool & Why it matters \\
\midrule
\texttt{collect} & gather a pipeline into one list \\
\texttt{partition} & split into matching/non-matching lists \\
\texttt{rename} & rename projected fields \\
\texttt{transpose} & pivot rows into columns \\
\texttt{zip}, \texttt{window}, \texttt{chunk} & sequence reshaping \\
\bottomrule
}

\note{\textbf{ToSh sweet spot:} load typed rows, project only the fields you care about, then summarize or regroup them without dropping down to fragile text parsing.}

\end{multicols*}
\newpage

\sheettitle{Collections And Pipelines}{Reshaping, folding, formatting, and advanced stream tools}{page 2 / reference}

\begin{multicols*}{3}

\section{Reshape Toolbox}
\tighttable{>{\ttfamily\raggedright\arraybackslash}p{0.27\linewidth}X}{
\toprule
Command & What it does \\
\midrule
\texttt{collect} & buffer the whole pipeline into one array \\
\texttt{chunk N} & fixed-size batches \\
\texttt{window N} & sliding windows over neighbors \\
\texttt{zip seq} & pair items with another sequence \\
\texttt{interleave seq} & alternate values from two sequences \\
\texttt{take-while} / \texttt{skip-while} & prefix/suffix slicing by predicate \\
\bottomrule
}

\begin{lstlisting}
echo 1 2 3 | collect
echo 1 2 3 4 | chunk 2
echo 1 2 3 4 | window 3
echo 1 2 3 | zip [10, 20, 30]
\end{lstlisting}

\section{Fold And Trace}
\begin{lstlisting}
echo 1 2 3 4 | reduce 0 func(acc, x) => ($acc + $x)
echo one two three | reduce "" { $acc + _.Substring(0, 1) }
echo 1 2 3 4 | scan 0 { $acc + _ }
\end{lstlisting}

\begin{itemize}
\item \texttt{reduce} returns one final accumulator.
\item \texttt{scan} yields every intermediate accumulator value.
\item On empty input, \texttt{reduce} returns the seed unchanged.
\end{itemize}

\section{Split And De-Dupe}
\begin{lstlisting}
var [evens, odds] = (echo 1 2 3 4 | partition { (_ % 2) == 0 })
echo alpha beta alpha gamma | distinct
echo alpha beta alpha gamma | frequencies
\end{lstlisting}

\section{Stopping And Slicing}
\begin{lstlisting}
echo 1 2 3 4 | take-while { _ < 3 }
echo 1 2 3 4 | skip-while { _ < 3 }
echo alpha beta gamma | first 2
echo alpha beta gamma | last 2
\end{lstlisting}

\columnbreak

\section{Projection Cleanup}
\begin{lstlisting}
ls -la |
get { Name, Size, Modified } |
rename Name FileName Size Bytes
\end{lstlisting}

\begin{lstlisting}
echo {| A = 1, B = 2 |} {| A = 3, B = 4 |} | transpose
\end{lstlisting}

\tighttable{>{\ttfamily\raggedright\arraybackslash}p{0.30\linewidth}X}{
\toprule
Command & Best for \\
\midrule
\texttt{rename} & clean projected field names \\
\texttt{transpose} & rows-to-columns pivoting \\
\texttt{pick}/\texttt{select} & aliases for \texttt{get} \\
\texttt{tee} & pass through while also capturing or writing \\
\bottomrule
}

\section{Structured Format Bridges}
\begin{lstlisting}
echo "{\"name\":\"toast\"}" | from json
curl https://example/api | from json | flatten
cat data.csv | from csv
\end{lstlisting}

\begin{lstlisting}
ls | to json
ls | to csv
ls | to json --compact
\end{lstlisting}

\note{\texttt{from} and \texttt{to} keep data as CLR objects until you explicitly flatten or stringify it, which is why downstream \texttt{get}, \texttt{where}, and \texttt{summarize} still work.}

\section{Text Templates}
\begin{lstlisting}
$rows | template "{{ Name }} => {{ Quantity }}"
$rows | template "{{ Category }} / {{ Name }}"
\end{lstlisting}

\section{Capture And Replay}
\begin{lstlisting}
ls | tee -v snapshot | first 3
$snapshot | get Name
ls | to json | tee report.json | count
\end{lstlisting}

\columnbreak

\section{Pipeline Recipes}
\begin{lstlisting}
func bigFiles(minSize: StorageSize) {
    ls -la | where _.Type == file |
    where _.Size >= $minSize |
    sort Size | reverse
}
\end{lstlisting}

\begin{lstlisting}
echo "[{\"name\":\"alpha\",\"score\":95},{\"name\":\"beta\",\"score\":87}]" |
from json | flatten | sort score | reverse | get { name, score }
\end{lstlisting}

\section{Grouping Recipes}
\begin{lstlisting}
ls | group-by Extension
ps | group-by func(p) => ($p.Name.Substring(0, 1))
echo 1 2 3 10 11 | group-while { _ < 10 }
\end{lstlisting}

\note{\texttt{group-by} groups by key across the whole stream, while \texttt{group-while} is about consecutive runs. Use the latter when input order already carries meaning.}

\section{Summary Patterns}
\begin{lstlisting}
df | summarize
df | summarize _.Used
seq 5 | summarize --sum --avg --min --max --count
ps | summarize --avg Memory --max Memory
\end{lstlisting}

\begin{itemize}
\item \texttt{sum}, \texttt{avg}, \texttt{min}, and \texttt{max} are great for one target.
\item \texttt{summarize} shines when you want multiple aggregates or structured column summaries back.
\end{itemize}

\section{Generated Sequences}
\begin{lstlisting}
iterate 1 { _ * 2 } | take 6
unfold 5 { if _ > 0 { [_, _ - 1] } else { null } }
converge "  toast  " { _.Trim() }
\end{lstlisting}

\note{\texttt{iterate} grows a stream from a seed, \texttt{unfold} emits value/state pairs until \texttt{null}, and \texttt{converge} keeps normalizing until two consecutive results are equal.}

\end{multicols*}
\end{document}
